\begin{figure}[ht]
\centering
\begin{tikzpicture}[
  vertex/.style = {circle, draw=engblack, thick, fill=englime!25,
                   inner sep=0pt, minimum size=5mm, font=\scriptsize},
  bvertex/.style = {circle, draw=engblack, thick, fill=engblue!15,
                    inner sep=0pt, minimum size=5mm, font=\scriptsize},
  edge/.style   = {thick, draw=engblack},
]

% --- Panel 1: planar / non-planar (left)
\begin{scope}[xshift=0cm]
\node[font=\small, anchor=south] at (1.5, 3.4) {(a) planar};
\node[vertex] (a1) at (0, 2)    {};
\node[vertex] (a2) at (3, 2)    {};
\node[vertex] (a3) at (1.5, 0)  {};
\node[vertex] (a4) at (1.5, 3)  {};
\draw[edge] (a1) -- (a2);
\draw[edge] (a1) -- (a3);
\draw[edge] (a2) -- (a3);
\draw[edge] (a1) -- (a4);
\draw[edge] (a2) -- (a4);
\end{scope}

% --- Panel 2: bipartite (middle)
\begin{scope}[xshift=5cm]
\node[font=\small, anchor=south] at (1.5, 3.4) {(b) bipartite $K_{3,3}$};
\foreach \i/\name in {0/u1, 1/u2, 2/u3} {
  \node[vertex] (\name) at (\i*1.5, 3) {};
}
\foreach \i/\name in {0/v1, 1/v2, 2/v3} {
  \node[bvertex] (\name) at (\i*1.5, 0) {};
}
\foreach \src in {u1, u2, u3}
  \foreach \dst in {v1, v2, v3}
    \draw[edge] (\src) -- (\dst);
\end{scope}

% --- Panel 3: tree (right)
\begin{scope}[xshift=10cm]
\node[font=\small, anchor=south] at (1.5, 3.4) {(c) tree};
\node[vertex] (r)  at (1.5, 3)   {};
\node[vertex] (c1) at (0.3, 2)   {};
\node[vertex] (c2) at (2.7, 2)   {};
\node[vertex] (g1) at (-0.3, 1)  {};
\node[vertex] (g2) at (0.9, 1)   {};
\node[vertex] (g3) at (2.1, 1)   {};
\node[vertex] (g4) at (3.3, 1)   {};
\node[vertex] (h1) at (-0.3, 0)  {};
\node[vertex] (h2) at (0.9, 0)   {};
\draw[edge] (r) -- (c1);
\draw[edge] (r) -- (c2);
\draw[edge] (c1) -- (g1);
\draw[edge] (c1) -- (g2);
\draw[edge] (c2) -- (g3);
\draw[edge] (c2) -- (g4);
\draw[edge] (g1) -- (h1);
\draw[edge] (g1) -- (h2);
\end{scope}

\end{tikzpicture}
\caption{Three graph families the engineer reads on sight.
\textbf{(a)} A planar graph drawable without edge crossings;
the complete graph $K_{4}$ here is the smallest non-trivial planar
graph that fills its faces. \textbf{(b)} The complete bipartite
$K_{3,3}$, the canonical non-planar graph and a Kuratowski
obstruction; bipartite structure makes it a stock model for
assignment problems. \textbf{(c)} A tree on nine vertices, eight
edges, height three; root at top.}
\label{fig:vol02:ch17:graph-families}
\end{figure}
